\documentclass[../differential_topology.tex]{subfiles}
\begin{document}
\section{Aula 14 - 29 de Outubro, 2026}
\subsection{Motivações}
\begin{itemize}
	\item Grau de Brouwer
\end{itemize}
\subsection{Grau de Brouwer}
\begin{def*}
	Sejam M e N orientáveis sem bordo com \(\mathrm{dim}_{}(M) = \mathrm{dim}_{}(N) = n\) e N conexa. Definimos:
	\begin{itemize}
		\item[i)] Se \(x\in M\) é um ponto regular de \(f:M\rightarrow N\), definimos o \textbf{sinal de \(\mathbf{f'(x)}\)} como
		      \[
			      \mathrm{sgn} f'(x) = \left\{\begin{array}{ll}
				      1, \quad & f'(x) \text{ preserva orientação} \\
				      -1,      & \text{caso contrário}
			      \end{array}\right.
		      \]
		\item[ii)] Para cada \(y\in N\) que seja valor regular de \(f:M\rightarrow N\), define-se o \textbf{grau de \(\mathbf{f}\) em \(\mathbf{y}\)} por
		      \[
			      \mathrm{deg}(f, y)=\sum\limits_{x\in f^{-1}(y)}^{}\mathrm{sgn}f'(x)
		      \]
	\end{itemize}
\end{def*}

\begin{prop*}
	No contexto anterior, quando M for uma variedade compacta, então \linebreak \(\mathrm{deg}(f, y)\) é um inteiro e é localmente constante.
\end{prop*}

\begin{lemma*}
	Suponha que \(M = \partial X\), onde X é uma variedade orientável com bordo. Se \(f:M\rightarrow N\) admite uma extensão \(F:X\rightarrow N\), então \(\mathrm{deg}(f) = 0\).
\end{lemma*}
\begin{lemma*}
	Se \(f\) e \(g\) são aplicações isotópicas, então para todo valor regular comum, vale
	\[
		\mathrm{deg}(f, y) = \mathrm{deg}(g, y).
	\]
\end{lemma*}

\begin{theorem*}
	O grau \(\mathrm{deg}(f, y)\) não depende do valor regular \(y\in N\) usado.
\end{theorem*}
\begin{proof*}
	Seja \(z\in N\) um valor regular de \(f\). Pelo \hyperlink{homogeneity_lemma}{\textit{Lema da Homogeneidade}}, existe uma aplicação \(h:N\rightarrow N\) tal que \(h(y) = z\) e que \(h\) é isotópica à identidade: \(\mathrm{id}_{N}\cong h\); assim, a composta \(h\circ f\) deverá ser isotópica à própria \(f\), tal que, pelo segundo lema,
	\[
		\mathrm{deg}(h\circ f, z) = \mathrm{deg}(f, z).
	\]

	Como um difeomorfismo só pode preservar ou reverter orientação, e a identidade é um que canonicamente preserva a orientação, é necessário que \(h\) e \(\mathrm{id}_{N}\) sejam conectados por um caminho de difeomorfismos, os quais todos preservam orientação.
	Logo, \(h\) preserva orientação, permitindo que calculemos o grau por
	\begin{align*}
		\mathrm{deg}(h\circ f, h(y)) = \sum\limits_{x\in (h\circ f)^{-1}(h(y))}^{} \mathrm{sgn}(h\circ f)'(x) & = \sum\limits_{x\in f^{-1}(y)}^{}\mathrm{sgn}(h'(f(x))\cdot f'(x)) \\
		                                                                                                      & =\sum\limits_{x\in f^{-1}(y)}^{}\mathrm{sgn}f'(x).
	\end{align*}
	Portanto, da relação
	\[
		\mathrm{deg}(f\circ h, z) = \mathrm{deg}(f, y) \Longleftrightarrow \mathrm{deg}(h\circ f, z) = \mathrm{deg}(f, z),
	\]
	segue a conclusão
	\[
		\mathrm{deg}(f, y) = \mathrm{deg}(f, z). \text{ \qedsymbol}
	\]

\end{proof*}
\begin{def*}
	O \textbf{grau de Brouwer de f} é, para algum \(y\in N\) valor regular de f,
	\[
		\mathrm{deg}(f) = \mathrm{deg}(f, y).\; \square
	\]
\end{def*}
\begin{theorem*}
	Sejam \(f, g:M\rightarrow N\); se \(g\) e \(f\) forem isotópicas, então seus graus são iguais.
\end{theorem*}

\begin{example}
	\begin{itemize}
		\item[1)] Se \(f:M\rightarrow N\) for a função constante, então
		      \[
			      \mathrm{deg}(f) = 0;
		      \]
		\item[2)] Se \(f:N\rightarrow N\) for a identidade, então
		      \[
			      \mathrm{deg}(\mathrm{id}) = \mathrm{deg}(f) = 1;
		      \]
		\item[3)] Dado um difeomorfismo qualquer \(f:M\rightarrow N\), vale
		      \[
			      \mathrm{deg}(f)  = \left\{\begin{array}{ll}
				      1, \quad  & f\text{ preserva orientação} \\
				      -1, \quad & f \text{ reverte orientação}
			      \end{array}\right.
		      \]
		\item[4)] Considere \(f:\mathbb{C}\rightarrow \mathbb{C}\) definida por \(f(z) = z^{k}\); então, \(f'(z) = k z^{k-1}\), mas olhando para isso na forma polar -- mais especificamente, o determinante disso -- nota-se que
		      \[
			      \mathrm{deg}f'(z) = \Vert kz^{k-1} \Vert^{2} > 0.
		      \]
		      Daí, qualquer número complexo não nulo é um valor regular, então basta tomarmos a pré-imagem de 1, \(f^{-1}(1) = \mathcal{U}_{n}\), onde \(\mathcal{U}_{n}\) é o conjunto das raízes da unidade
		      \[
			      \mathcal{U}_{n} = \{z: z^{n-1} = 0\},
		      \]
		      tal que
		      \[
			      \mathrm{deg}(f, 1) = \mathrm{deg}(f) =\sum\limits_{z\in f^{-1}(1)}^{}1 = k,
		      \]
		      pois para \(z^{k}\) teremos k raízes da unidade;
		\item[5)] Seja \(R_{i}:\mathbb{S}^{n-1}\rightarrow \mathbb{S}^{n-1}\) a reflexão na i-ésima coordenada, \textit{i.e.},
		      \[
			      R_{i}(x_1, \dotsc , x_{n}) = (x_1, \dotsc ,-x_{i}, \dotsc , x_{n}).
		      \]
		      Neste caso, \(\mathrm{deg}(R_{i}) = -1\), e isso é um resultado muito importante (mas falta um teorema ainda pra entender como), pois permite que analisemos o grau da aplicação antipodal \(A:\mathbb{S}^{n-1}\rightarrow \mathbb{S}^{n-1}\), dada pela composição de todas as \(R_{i}\)'s:
		      \[
			      A \coloneqq R_1\circ \dotsc R_{n},
		      \]
		      ou seja, se soubermos como calcular o grau de uma composição de funções, podemos encontrar o grau da aplicação antipodal (Spoiler: \(\mathrm{deg}(A) = (-1)^{n}\)).
	\end{itemize}
\end{example}

\end{document}
